\documentclass{amsart}

% 数学相关宏包
\usepackage{amsmath,mathtools,amsthm,amssymb}
\usepackage{bbm}
\usepackage{mathrsfs}
\usepackage{centernot}

% 图形和绘图
\usepackage{graphicx,float}
\usepackage{tikz,tikz-cd}
\usepackage{pgfplots}
\usepackage[framemethod=TikZ]{mdframed}
\usetikzlibrary{decorations.pathmorphing,positioning,arrows,automata,calc,shapes,patterns}
\pgfplotsset{compat=1.18}
\usepgfplotslibrary{statistics}

% 表格和列表
\usepackage{booktabs,multirow,colortbl}
\usepackage{enumitem}
\setlist[itemize,enumerate,description]{noitemsep}
\setlistdepth{9}

% 算法
\usepackage{algorithm,algorithmic}

% 字体和颜色
\usepackage{xcolor,listings}
\usepackage{xeCJK}
\setCJKmainfont{Noto Serif CJK SC}

% 页面和链接
\usepackage{fancyhdr,caption}
\usepackage{hyperref}
\usepackage{cleveref}

% 工具宏包
\usepackage{etoolbox,letltxmacro,docmute}
\usepackage{gensymb}

% 自定义设置
\renewcommand{\arraystretch}{1.5}
\let\originalmathbb\mathbb
\renewcommand{\mathbb}[1]{%
    \ifstrequal{#1}{1}{\mathbbm{#1}}{\originalmathbb{#1}}%
}

% 颜色定义
\definecolor{lightblue}{RGB}{173,216,230}
\definecolor{lightgreen}{RGB}{144,238,144}
\definecolor{lightyellow}{RGB}{255,255,224}
\definecolor{lightgray}{RGB}{211,211,211}

% 超链接设置
\hypersetup{
    colorlinks=true,
    linkcolor=red,
    filecolor=magenta,
    urlcolor=cyan,
    pdftitle={\@title},
    pdfauthor={\@author},
    pdfsubject={数学笔记},
    pdfkeywords={数学},
    bookmarksnumbered=true,
    bookmarksopen=true,
    bookmarksopenlevel=6,
    pdfstartview=FitH,
    pdfpagemode=UseOutlines,
    pdfborder={0 0 0}
}

% 定理环境
\theoremstyle{plain}
\newtheorem{theorem}{Theorem}[section]
\newtheorem{lemma}[theorem]{Lemma}
\newtheorem{corollary}[theorem]{Corollary}
\newtheorem{proposition}[theorem]{Proposition}

\theoremstyle{definition}
\newtheorem{definition}[theorem]{Definition}
\newtheorem{example}[theorem]{Example}
\newtheorem{exercise}[theorem]{Exercise}
\newtheorem{problem}[theorem]{Problem}

\theoremstyle{remark}
\newtheorem{remark}[theorem]{Remark}
\newtheorem{note}[theorem]{Note}
\newtheorem{observation}[theorem]{Observation}

% 解答环境
\newtheorem*{solution}{Solution}
\newtheorem*{proof*}{Proof}

% 图片路径
\graphicspath{{F:/iCloudDrive/iCloud~md~obsidian/2025-summer/figures/}}

\author{乐绎华, Yue Yihua}
\address{中山大学, Sun Yat-sen University}
\email{yueyh@mail2.sysu.edu.cn}
\date{\today}
\title{Mathematical Notes}

\begin{document}

\maketitle
\tableofcontents

\section{Why study Grassmannian?}

\begin{definition}[Grassmannian $\mathrm{Gr}(k, n)$]
The \textbf{Grassmannian} $\mathrm{Gr}(k, n)$ is the set of all $k$-dimensional linear subspaces of an $n$-dimensional vector space $\mathbb{F}^n$ (where $\mathbb{F}$ is usually $\mathbb{R}$ or $\mathbb{C}$).
Formally:
\[
\mathrm{Gr}(k, n) = \{ \text{\(k\)-dimensional subspaces of } \mathbb{F}^n \}.
\]	\begin{itemize}
		\item Example: $\mathrm{Gr}(1, n)$ is the \textbf{projective space} $\mathbb{P}^{n-1}$, since a 1D subspace is just a line through the origin.
		\item Example: $\mathrm{Gr}(n-1, n)$ can be thought of as the set of hyperplanes through the origin.
	\end{itemize}
The Grassmannian is not just a set — it is a \textbf{smooth manifold} (and even an \textbf{algebraic variety} when $\mathbb{F} = \mathbb{C}$).
Its \textbf{dimension} is:
\[
\dim_{\mathbb{F}} \mathrm{Gr}(k, n) = k(n-k).
\]
\end{definition}
Imagine $\mathbb{R}^n$. You can look at:

\begin{itemize}
	\item Points ($k=0$)
	\item Lines through the origin ($k=1$)
	\item Planes through the origin ($k=2$)
	\item etc.
\end{itemize}

The Grassmannian $\mathrm{Gr}(k, n)$ is like a "space whose points are subspaces" — you can \textbf{move smoothly from one subspace to another}. This lets us treat "families of subspaces" as geometric objects themselves.

A key idea:
The Grassmannian is \textbf{homogeneous}:
\[
\mathrm{Gr}(k, n) \cong \frac{\mathrm{GL}(n, \mathbb{F})}{P},
\]
where $P$ is a \textbf{parabolic subgroup} that fixes a chosen $k$-dimensional subspace. This shows that all $k$-planes look "the same" up to a linear transformation.

When $\mathbb{F} = \mathbb{C}$, $\mathrm{Gr}(k, n)$ can be embedded into projective space via \textbf{Plücker coordinates}:
\[
V \in \mathrm{Gr}(k, n) \quad \mapsto \quad v_1 \wedge \cdots \wedge v_k \in \mathbb{P}\big(\Lambda^k \mathbb{C}^n\big),
\]
where $\{v_i\}$ is a basis for $V$.

The image is cut out by \textbf{quadratic Plücker relations}, making the Grassmannian a \textbf{projective algebraic variety}.

This algebraic viewpoint is key in \textbf{Schubert calculus}, where we study how subspaces intersect.

\subsection{\textbf{Linear Algebra \& Geometry}}

\begin{itemize}
	\item \textbf{Parametrizing solutions}: If you have a family of subspaces (e.g., solution spaces to a system of equations), they can be described as points in a Grassmannian.
	\item \textbf{Change of basis}: Studying all $k$-subspaces is like studying all possible column spaces of $n \times k$ matrices up to linear transformations.
\end{itemize}

\subsection{\textbf{Algebraic Geometry}}

\begin{itemize}
	\item Grassmannians are \textbf{building blocks}: Many moduli spaces can be constructed from them.
	\item \textbf{Schubert varieties} inside Grassmannians help compute intersection numbers (enumerative geometry problems like “how many lines in $\mathbb{P}^3$ meet four given lines?”).
\end{itemize}

\subsection{\textbf{Topology \& Characteristic Classes}}

\begin{itemize}
	\item Grassmannians are classifying spaces for vector bundles:
\end{itemize}
\[
\mathrm{Gr}_k(\mathbb{R}^\infty) \quad\text{classifies rank-\(k\) real vector bundles}.
\]
\begin{itemize}
	\item Their cohomology ring is generated by \textbf{Chern/Schubert classes}, essential in topology and geometry.
\end{itemize}

\subsection{\textbf{Representation Theory}}

\begin{itemize}
	\item As homogeneous spaces $G/P$, Grassmannians are examples of \textbf{flag varieties}, and their geometry reflects the structure of Lie groups and Lie algebras.
\end{itemize}

\subsection{\textbf{Physics}}

\begin{itemize}
	\item In \textbf{quantum field theory} and \textbf{string theory}, Grassmannians appear in:
	\begin{itemize}
		\item \textbf{Gauge theory}: Moduli of certain field configurations are often Grassmannians or their generalizations.
		\item \textbf{Twistor theory}: The "momentum space" in some scattering amplitude formulations is a Grassmannian.
	\end{itemize}
\end{itemize}

\subsection{\textbf{Applied Mathematics}}

\begin{itemize}
	\item \textbf{Statistics \& Data Science}: Subspaces model principal components in PCA; the set of all $k$-dimensional subspaces is the search space in \textbf{subspace methods}.
	\item \textbf{Signal Processing}: In array processing and MIMO systems, signals lie in subspaces, and the Grassmannian provides a natural search domain.
\end{itemize}

\begin{table}[h]
	\centering
	\begin{tabular}{|c|c|}
		\hline
		Aspect & Why it matters \\
		\hline
		\textbf{Geometry} & Smooth, compact, homogeneous manifold parametrizing subspaces \\
		\hline
		\textbf{Algebra} & Projective variety with explicit Plücker equations \\
		\hline
		\textbf{Topology} & Classifying space for vector bundles \\
		\hline
		\textbf{Combinatorics} & Schubert calculus counts geometric configurations \\
		\hline
		\textbf{Applications} & From quantum physics to machine learning \\
		\hline
	\end{tabular}
\end{table}
If you like, I can \textbf{follow up} by explaining \textbf{Schubert varieties inside the Grassmannian}, which are one of the main computational tools for studying it — and they connect directly to the \textit{applications in enumerative geometry}. That’s where the Grassmannian really starts to \textit{do things}.


\end{document}